\chapter{Synchronous Model}

\begin{center}
{\Large\bfseries Synchronous Model}\\[1pt]
{\small Exam cheat-sheet: \emph{execution traces, reaching a target state, extended state machines}}
\end{center}
\hrule
\vspace{4pt}

Exam Exercise~1 on synchronous components is almost always the same three tasks
on a given component:
\textbf{(a)} fill in an execution trace, \textbf{(b)} find inputs reaching a
target state, \textbf{(c)} redraw the component as an extended state machine.

%==================================================================
\section{Notation cheat-sheet}

\begin{itemize}
  \item \textbf{Component} $=(I,O,S,\mathrm{Init},\mathrm{React})$: input vars,
        output vars, state vars, initialisation, reaction code run \emph{once per round}.
  \item \textbf{Reaction:} $\;s \xrightarrow{\,i/o\,} t\;$ means: in old state $s$,
        with input $i$, produce output $o$ and move to new state $t$.
  \item \textbf{State} is written as the tuple of state variables in declaration
        order, e.g.\ $(\mathit{mode},x)$. Inputs are written as the tuple of input
        vars; output after the ``$/$''.
  \item \textbf{Events:} \texttt{event} $\in\{\bot,\top\}$;
        \texttt{event(nat)} $\in\{\bot,0,1,2,\dots\}$.
        $\bot$ = absent, $\top$ = present.
  \item \texttt{x?} tests ``$x$ present''. \texttt{x!v} means \texttt{x:=v}
        (and \texttt{x!} means \texttt{x:=present}).
        \textbf{Unassigned output event $\Rightarrow$ absent ($\bot$) by default.}
  \item \textbf{Combinational} = \emph{no} state variables (e.g.\ \texttt{out:=in1-in2}).
        \textbf{Sequential} = has state variables (e.g.\ Delay).
  \item \textbf{Deterministic} = single initial state \emph{and} for every $(s,i)$
        exactly one $(o,t)$. (\texttt{choose}/nondeterminism breaks this.)
\end{itemize}

%==================================================================
\section{(a) Fill in an execution trace}

\textbf{Recipe (one round at a time):}
\begin{enumerate}
  \item Write the current state tuple and plug the given input values in.
  \item Run the update code \emph{top to bottom}, line by line, keeping a scratch
        value for each variable.
  \item Read off the \textbf{output} ($o$) and the \textbf{new state} ($t$).
  \item Carry $t$ forward as the start state of the next round; repeat.
\end{enumerate}
\begin{keybox}
\textbf{Trap:} if the code first flips the \texttt{mode} and \emph{then} uses
\texttt{mode}, the later lines use the \emph{new} mode. Always update in code order.
\end{keybox}

\textbf{Worked (AddOrMultiply, state $(\mathit{mode},x)$, input $(\mathit{in},\mathit{flip})$):}
code is ``\texttt{if flip? then flip mode}; then if \texttt{add}: \texttt{x:=x+in}
else \texttt{x:=x*in}; \texttt{out:=x}''.
\[
(\mathrm{add},0)
\xrightarrow{(2,\bot)/2}(\mathrm{add},2)
\xrightarrow{(3,\top)/6}(\mathrm{mul},6)
\xrightarrow{(1,\top)/7}(\mathrm{add},7)
\xrightarrow{(3,\bot)/10}(\mathrm{add},10)
\]
Round~2 check: \texttt{flip} present $\Rightarrow$ mode becomes \texttt{mul};
\emph{then} \texttt{x:=x*in}$=2\cdot3=6$; \texttt{out}$=6$.

%==================================================================
\section{(b) Find inputs reaching a target state}

\textbf{Recipe:}
\begin{enumerate}
  \item Note the target tuple, e.g.\ $(\mathrm{mul},42)$, and the initial state.
  \item Work out the \emph{effect of one input} on each variable (what makes the
        mode switch, what arithmetic is applied).
  \item Find the \textbf{minimum length}: if you must reach a mode you cannot start
        in, you need at least one switching input, so often length $\ge 2$.
  \item Construct concrete inputs that produce the target value; remember to set
        each event to $\bot$ or $\top$ as required.
\end{enumerate}

\textbf{Worked:} reach $(\mathrm{mul},42)$ from $(\mathrm{add},0)$.
First input $(a,\bot)$ stays in \texttt{add} and sets $x:=0+a=a$; second input
$(b,\top)$ flips to \texttt{mul} then sets $x:=a\cdot b$. Need $a\cdot b=42$, e.g.
\[
(\mathrm{add},0)\xrightarrow{(42,\bot)/42}(\mathrm{add},42)
\xrightarrow{(1,\top)/42}(\mathrm{mul},42).
\]
Any factor pair works: $(42,\bot)(1,\top)$, $(21,\bot)(2,\top)$, $(7,\bot)(6,\top)$, \dots

%==================================================================
\section{(c) Redraw as an extended state machine}

\textbf{Recipe:}
\begin{enumerate}
  \item The named state variable (e.g.\ \texttt{mode}) becomes the
        \textbf{modes/bubbles}: one bubble per possible value.
  \item Every \emph{other} state variable stays a \textbf{data variable},
        declared with its initial value on the start-arrow
        (e.g.\ \texttt{int x:=0}). Put the start arrow into the initial mode.
  \item For each mode and each relevant case of the inputs, draw a transition.
        Find its \textbf{target mode} by simulating the mode-switch logic, and label it
        \texttt{guard? $\to$ updates; out:=\dots}.
        Use a self-loop when the mode does not change.
  \item Updates on a switching transition use the behaviour of the
        \emph{destination} mode (same trap as in (a)).
\end{enumerate}

\textbf{Worked (AddOrMultiply)}, modes \texttt{add} and \texttt{mul}, init \texttt{add},
data \texttt{int x:=0}:
\begin{itemize}
  \item \texttt{add} $\xrightarrow{\;\neg\texttt{flip?}\;\to\;x:=x+in;\ out:=x\;}$ \texttt{add}\quad(self-loop)
  \item \texttt{add} $\xrightarrow{\;\texttt{flip?}\;\to\;x:=x*in;\ out:=x\;}$ \texttt{mul}
  \item \texttt{mul} $\xrightarrow{\;\neg\texttt{flip?}\;\to\;x:=x*in;\ out:=x\;}$ \texttt{mul}\quad(self-loop)
  \item \texttt{mul} $\xrightarrow{\;\texttt{flip?}\;\to\;x:=x+in;\ out:=x\;}$ \texttt{add}
\end{itemize}
(The two switching edges cross: \texttt{add}$\to$\texttt{mul} \emph{multiplies},
\texttt{mul}$\to$\texttt{add} \emph{adds}, because the arithmetic uses the new mode.)

%==================================================================
\section{Quick facts for short-answer questions}

\begin{itemize}
  \item ``Is it combinational?'' $\Rightarrow$ Yes iff it has no state variables.
  \item ``Is it deterministic?'' $\Rightarrow$ Yes iff one initial state and a unique
        $(o,t)$ for every $(s,i)$; a \texttt{choose} makes it nondeterministic.
  \item \textbf{Synchrony hypothesis:} update code is instantaneous; outputs produced
        and inputs read in the \emph{same} round; all composed components fire
        simultaneously each round.
  \item Event-triggered component reacts only in rounds where its trigger event is present.
\end{itemize}

\vspace{4pt}
\hrule
\vspace{2pt}
{\footnotesize\textbf{Exam triage:} \emph{trace?} run code top-to-bottom each round; new mode takes effect immediately. \emph{reach target?} compute mode-switch cost, pick factor pair for numeric target. \emph{ESM?} mode var $\to$ bubbles; other state vars $\to$ data on start-arrow; one edge per (mode, guard) case; switching edges use destination-mode arithmetic. \emph{combinational?} no state vars. \emph{deterministic?} unique $(o,t)$ for every $(s,i)$.}
